% ============================================================
% MAPA STATUSU TEORII TGP — Dodatek A.1 (sesja v33)
% ============================================================
\section{Mapa statusu twierdze\'n TGP}\label{sec:status-map}

\begin{center}
\fbox{\parbox{0.9\linewidth}{%
\textbf{Nota kanoniczna (2026-04-07).}
Warto\'sci $g_0 = 1{,}24$ w~tym dokumencie odpowiadaj\k{a} Formulacji~B ($f(g) = 1+4\ln g$).
Warto\'s\'c kanoniczna (Formulacja~A, $K(g) = g^4$): $g_0^e = 0{,}86941$.
Algebra Koide'a i~r\'ownania F1--F12 s\k{a} niezale\.zne od tego wyboru.}}
\end{center}

\begin{remark}[Cel i~konwencja]\label{rem:status-map-intro}
Niniejsza sekcja pe\l{}ni rol\k{e} \textbf{indeksu logicznego} teorii.
Ka\.zde zdanie twierdzeniowe TGP jest opatrzone jednym z~sze\'sciu
etykiet statusu (uw.~\ref{rem:status-levels}):
\statuslabel{Aksjomat},
\statuslabel{Twierdzenie},
\statuslabel{Propozycja},
\statuslabel{Hipoteza},
\statuslabel{Szkic},
\statuslabel{Program}.
Tabela poni\.zej zbiera wszystkie kluczowe wyniki
i~pozwala czytelnikowi natychmiast oceni\'c,
na jakim gruncie stoj\k{a}.
\end{remark}

% -------------------------------------------------------
\subsection*{Warstwa~0: Substrat i~operator~$\mathcal{D}$}
% -------------------------------------------------------

\begin{center}\footnotesize
\begin{tabular}{@{}p{5.2cm}lp{4.4cm}p{2.6cm}@{}}
\toprule
\textbf{Zdanie} & \textbf{Status} & \textbf{Warunki / za\l{}o\.zenia}
  & \textbf{Odsylacz} \\
\midrule

Substrat $\Gamma = (V,E)$ z~symetri\k{a} $\mathbb{Z}_2$
i~hamiltonianem $H_\Gamma$ istnieje
  & \statuslabel{Aksjomat}
  & (A1) --- wej\'scie teorii
  & aks.~A1 \\[3pt]

$K(0) = 0$: brak przestrzeni przy $\Phi = 0$
  & \statuslabel{Twierdzenie}
  & $K(\varphi) = K_{\rm geo}\varphi^4 \Rightarrow K(0) = 0$; wyprowadzenie z~substratu
  & stw.~\ref{prop:K0-from-substrate} \\[3pt]

$\alpha = 2$: operator $\mathcal{D} = \nabla^2\Phi/\Phi + 2(\nabla\Phi)^2/\Phi^2$
  & \statuslabel{Twierdzenie}
  & $K_{ij} = J(\varphi_i\varphi_j)^2$, klasa $K(\varphi) = \varphi^n$
  & tw.~\ref{prop:substrate-action} \\[3pt]

$\beta = \gamma$: warunek pr\'o\.zniowy $U'(1) = 0$
  & \statuslabel{Twierdzenie}
  & $U(\psi)$ minimalne w~$\psi=1$
  & stw.~\ref{prop:vacuum-condition} \\[3pt]

$d = 3$: trzy wymiary przestrzenne substratu
  & \statuslabel{Twierdzenie}
  & optymalno\'s\'c perkolacyjna sieci
  & stw.~\ref{prop:wymiar} \\[3pt]

Jednoznaczno\'s\'c operatora $\mathcal{D}$ w~klasie
drugiego rz\k{e}du bez zewn\k{e}trznej metryki
  & \statuslabel{Twierdzenie}
  & warunki (1)--(4) z~\S\ref{ssec:rownanko}
  & tw.~\ref{thm:D-uniqueness} \\[3pt]

Pr\'o\.znia z~\'zr\'od\l{}em: $\mathcal{D}[\Phi](x) = -q\rho_{\rm mat}
  - \Phi_0 W(\psi)$, $W(1) = \gamma/3 \neq 0$
  & \statuslabel{Twierdzenie}
  & $\beta = \gamma$, dzia\l{}anie~\eqref{eq:action}
  & tw.~\ref{thm:vacuum-source} \\[3pt]

$m_{\rm sp} = \sqrt{\gamma}$: masa bozonu przestrzennego
  & \statuslabel{Twierdzenie}
  & N0-5 $\Rightarrow$ N0-6; linearyzacja r\'ownania pola (sesja v33)
  & stw.~\ref{prop:N0-6-from-N0-5} \\[3pt]

Mapa rodziny sprz\k{e}\.ze\'n $K(\varphi) = \varphi^n$:
$K = \varphi^4$ jest minimalne dla naszego wszech\'swiata
  & \statuslabel{Propozycja}
  & klasa $K(\varphi) = \varphi^n$, obs. $c_T = c_0$
  & uw.~\ref{rem:coupling-map} \\

\bottomrule
\end{tabular}
\end{center}

% -------------------------------------------------------
\subsection*{Warstwa~1b: Dynamiczne sta\l{}e fundamentalne}
% -------------------------------------------------------

\begin{center}\footnotesize
\begin{tabular}{@{}p{5.2cm}lp{4.4cm}p{2.6cm}@{}}
\toprule
\textbf{Zdanie} & \textbf{Status} & \textbf{Warunki / za\l{}o\.zenia}
  & \textbf{Odsylacz} \\
\midrule

$c(\Phi) = c_0(\Phi_0/\Phi)^{1/2}$
  & \statuslabel{Propozycja}
  & metryka pomostowa + s\l{}abe pole
  & stw.~\ref{prop:c-from-metric} \\[3pt]

$\hbar(\Phi) = \hbar_0(\Phi_0/\Phi)^{1/2}$
  & \statuslabel{Propozycja}
  & $c/\hbar = \mathrm{const}$ (warunek~W3)
  & stw.~\ref{prop:G-from-lP}, tw.~\ref{thm:exponents} \\[3pt]

$G(\Phi) = G_0\,\Phi_0/\Phi$
  & \statuslabel{Twierdzenie}
  & $\ell_P = \mathrm{const}$ (tw.~\ref{thm:lP})
  & stw.~\ref{prop:G-from-lP} \\[3pt]

$\ell_P = \sqrt{\hbar_0 G_0/c_0^3} = \mathrm{const}$
  & \statuslabel{Twierdzenie}
  & kombinacja wyk\l{}adnik\'ow: $b+g-3a = 0$
  & tw.~\ref{thm:lP} \\[3pt]

Jednoznaczno\'s\'c wyk\l{}adnik\'ow $(a,b,g) = (1/2,1/2,1)$
  & \statuslabel{Twierdzenie}
  & warunki W1--W3; uk\l{}ad 3 r\'owna\'n, rz\k{a}d~3
  & tw.~\ref{thm:exponents} \\[3pt]

Zasada nieoznaczono\'sci $\Delta x \cdot \Delta p \geq \hbar(\Phi)$
  & \statuslabel{Propozycja}
  & model energii pola rdzeniowego; $f = 1$ jawny
  & stw.~\ref{prop:uncertainty-tgp} \\[3pt]

$a_\Gamma \cdot \Phi_0 = 1$: redukcja $N_{\rm param} = 2 \to 1$
  & \statuslabel{Hipoteza}
  & DESI DR2: $a_\Gamma \Phi_0 = 1{,}005$ ($1{,}03\sigma$);
    Planck: $0{,}987$ ($1{,}22\sigma$);
    \texttt{ex110} (10/10 PASS)
  & R3; p75 \\

\bottomrule
\end{tabular}
\end{center}

% -------------------------------------------------------
\subsection*{Warstwa~2: Metryka efektywna i~grawitacja}
% -------------------------------------------------------

\begin{center}\footnotesize
\begin{tabular}{@{}p{5.2cm}lp{4.4cm}p{2.6cm}@{}}
\toprule
\textbf{Zdanie} & \textbf{Status} & \textbf{Warunki / za\l{}o\.zenia}
  & \textbf{Odsylacz} \\
\midrule

Forma metryki: diagonalna i~izotropowa
($g_{\mu\nu} \propto \delta_{\mu\nu}$)
  & \statuslabel{Twierdzenie}
  & symetria $\mathbb{Z}_2$ + izotropia substratu (sesja v33)
  & tw.~\ref{thm:metric-form-uniqueness} \\[3pt]

Metryka efektywna $g_{\mu\nu} = f(\Phi)\,\eta_{\mu\nu}$
(skalarna funkcja $f$)
  & \statuslabel{Propozycja}
  & tw.~\ref{thm:metric-form-uniqueness} + aks.~A3 + A6;
    $f(\Phi) = \Phi/\Phi_0$ z~norm.~czasu
  & stw.~\ref{prop:g00-from-axioms} \\[3pt]

Granica Newtona $V_N \propto -1/r$
  & \statuslabel{Twierdzenie}
  & metryka z~hipotezy pomostowej
  & tw.~\ref{thm:newton} \\[3pt]

PPN: $\gamma_{\rm PPN} = \beta_{\rm PPN} = 1$
  & \statuslabel{Twierdzenie}
  & metryka konformalna (izotropowa)
  & stw.~\ref{rem:PPN-Will} \\[3pt]

$c_{\rm GW} = c_0$ (pr\k{e}dko\'s\'c fal GW)
  & \statuslabel{Twierdzenie}
  & jednorodne t\l{}o $\Phi$, metryka pomostowa
  & stw.~\ref{prop:cT} \\[3pt]

Emergencja r\'owna\'n Einsteina do $O(U^2)$
  & \statuslabel{Propozycja}
  & metryka pomostowa + rozw.\ perturbacyjne
  & tw.~\ref{thm:einstein-emergence} \\[3pt]

$M_* = m_P$: skala Plancka z~substratu
  & \statuslabel{Propozycja}
  & wymiarowanie $\Phi_0 + \ell_P$; brak mikro-derywacji
  & stw.~\ref{prop:Mstar-from-substrate} \\[3pt]

$\sigma_{ab}$: mody tensorowe z~tensora
napr\k{e}\.ze\'n substratu
  & \statuslabel{Hipoteza}
  & rozszerzenie A1; nie wynika z~samego $\Phi$
  & \S\ref{ssec:tensor-substrate} \\[3pt]

Warunek antypodyczny $f \cdot h = 1$ z~bud\.zetu substratu
  & \statuslabel{Propozycja [AN+POST]}
  & $\mathcal{B} = N_B \cdot s_0 = \text{const}$; podzia\l{} czas/przestrze\'n;
    za\l{}o\.zenie multiplikatywne
  & stw.~\ref{prop:antipodal-from-budget} \\[3pt]

Metryka z~substratu: $g_{tt} = -c_0^2/\psi$, jedyno\'s\'c
  & \statuslabel{Twierdzenie}
  & (i) $g_{ij}$; (ii) $fh{=}1$; (iii) PPN; (iv) $\ell_P{=}\text{const}$
  & tw.~\ref{thm:metric-from-substrate-full} \\[3pt]

Zunifikowana akcja $S_{\rm TGP}$; $\sqrt{-g} = c_0\psi$
  & \statuslabel{Propozycja}
  & metryka $\to$ $\det g = -c_0^2\psi^2$
  & hip.~\ref{hyp:unified-action} \\[3pt]

$\kappa = 3/(4\Phi_0) \approx 0{,}030$ (sprz\k{e}\.zenie
materia--pole)
  & \statuslabel{Twierdzenie}
  & $\delta S/\delta\Phi = 0$ w~granicy FRW
  & stw.~\ref{prop:kappa-corrected} \\[3pt]

Brak ghost\'ow w~sektorze solitonowym
  & \statuslabel{Twierdzenie}
  & $K_{\rm sub}(g) = g^2 > 0\ \forall g > 0$;
    $f(g)$ to przybli\.zenie
  & tw.~\ref{thm:ghost-free-soliton} \\

\bottomrule
\end{tabular}
\end{center}

% -------------------------------------------------------
\subsection*{Kosmologia i~perturbacje}
% -------------------------------------------------------

\begin{center}\footnotesize
\begin{tabular}{@{}p{5.2cm}lp{4.4cm}p{2.6cm}@{}}
\toprule
\textbf{Zdanie} & \textbf{Status} & \textbf{Warunki / za\l{}o\.zenia}
  & \textbf{Odsylacz} \\
\midrule

$\Lambda_{\rm eff} \approx \gamma/12$: stała kosmologiczna
  & \statuslabel{Twierdzenie}
  & $W(1) = \gamma/3$, $\psi \approx 1$
  & stw.~\ref{prop:Lambda-eff} \\[3pt]

Równanie Friedmanna TGP~\eqref{eq:friedmann-modified}
  & \statuslabel{Propozycja}
  & metryka pomostowa; wyprowadzenie warunkowe
  & stw.~\ref{prop:cosmo} \\[3pt]

Pe\l{}ne r\'ownanie TGP w~FRW z~gradientami
  & \statuslabel{Twierdzenie}
  & A1 + metryka pomostowa + dzia\l{}anie~\eqref{eq:action}
  & stw.~\ref{prop:TGP-FRW-full} \\[3pt]

Zmienna MS: $v_k = a\psi_{\rm bg}^2\,\delta\psi_k$
  & \statuslabel{Twierdzenie}
  & pochodne dzia\l{}ania, $K(\psi) = \psi^4$
  & stw.~\ref{prop:MS-TGP} \\[3pt]

Spektrum pierwotne $\mathcal{P}_s(k) = H^2/(4\pi^2 c_0\psi_0^4)$
  & \statuslabel{Propozycja}
  & de~Sitter + zmienne~MS; brak korekty $\epsilon$
  & stw.~\ref{cor:TGP-power-spectrum} \\[3pt]

Tilt $(n_s - 1)_{\rm TGP} = -4\varepsilon_H - 4\varepsilon_\psi$
  & \statuslabel{Propozycja}
  & wzorzec $z''/z$ + numeryka \texttt{ex54}; w~oknie Plancka
  & r.~\eqref{eq:ns-TGP}, \texttt{ex54} \\[3pt]

$\psi_{\rm ini} = 7/6$ (atraktor GL + BBN)
  & \statuslabel{Propozycja}
  & $W(7/6)=0$ + $\Delta G/G=0$ + GL;
    \texttt{ex86} 9/9~PASS (sesja v35)
  & stw.~\ref{prop:psi-ini-derived} \\

\bottomrule
\end{tabular}
\end{center}

% -------------------------------------------------------
\subsection*{Sektor materii (poziom~3)}
% -------------------------------------------------------

\begin{center}\footnotesize
\begin{tabular}{@{}p{5.2cm}lp{4.4cm}p{2.6cm}@{}}
\toprule
\textbf{Zdanie} & \textbf{Status} & \textbf{Warunki / za\l{}o\.zenia}
  & \textbf{Odsylacz} \\
\midrule

$N_{\rm gen} = 3$: dok\l{}adnie trzy generacje
  & \statuslabel{Twierdzenie}
  & WKB $n = 0,1,2$ przy $a_\Gamma$; $N_{\rm gen} < 4$
  & stw.~\ref{prop:no-4th-generation} \\[3pt]

Statystyki Fermiego kinku: $\Psi(2,1) = -\Psi(1,2)$
  & \statuslabel{Twierdzenie}
  & $\pi_1(\mathcal{C}_{\rm sol}) = \mathbb{Z}_2$; tw.~Laidlawa--DeWitta;
    formalny dow\'od w~dod.~E$'$
  & tw.~\ref{thm:pi1-Z2}, wn.~\ref{cor:fermions} \\[3pt]

Formu\l{}a Koidego $Q = 2/3$
  & \statuslabel{Propozycja}
  & $\alpha_K$ dopasowane; pred. wzgl\k{e}dnych mas
  & \S\ref{ssec:koide-hierarchy} \\[3pt]

$\alpha_K \approx 8{,}56$ z~warunku bifurkacji $a_c(\alpha_K) = a_\Gamma$
  & \statuslabel{Hipoteza}
  & OP-3 otwarty; \'scie\.zki 1--8 zamkni\k{e}te negatywnie
  & uw.~\ref{rem:parsimony}; p104--p112 \\[3pt]

\'{S}cie\.zka~9: $\varphi$-FP, $r_{21} = 206{,}77$ (odch.~$0{,}0001\%$)
  & \statuslabel{Twierdzenie}
  & \texttt{ex106} (14/14); $g_0^* = 1{,}24915$;
    $g_0^\mu = \varphi g_0^* = 2{,}02117$; zero param.\ wolnych
  & twierdzenie~\ref{thm:J2-FP}; dod.~\ref{app:J2-formalizacja} \\[3pt]

$M_{\rm TGP} = c_M A_{\rm tail}^4 + \mathcal{O}(A^6)$: tryb zerowy
  & \statuslabel{Twierdzenie [AN+NUM]}
  & $\mathcal{E}^{(2)}{=}0$ [AN]; $\mathcal{E}^{(3)}{=}0$ [\textbf{NUM}]
    (dow\'od analityczny otwarty, uw.~\ref{rem:R-E3-status});
    \texttt{ex142}: 12/12 PASS
  & tw.~\ref{thm:R-A4}; dod.~\ref{app:zero-mode-A4} \\[3pt]

Warunek kwantowania $B_{\rm tail}(g_0^e) = 0$ (elektron)
  & \statuslabel{Propozycja}
  & $g_0^e = 1{,}2301$ (0{,}8\%); \texttt{ex58}
  & stw.~\ref{prop:J-quantization}; dod.~\ref{app:ogon-masy} \\[3pt]

Z\l{}ota proporcja $g_0^\mu = \varphi \cdot g_0^e$: $\varphi$-FP
  & \statuslabel{Twierdzenie}
  & wynika z~$\varphi$-FP (twierdzenie~\ref{thm:J2-FP});
    nie za\l{}o\.zenie, lecz wynik dynamiczny
  & def.~\ref{def:J2-phi-FP}; dod.~\ref{app:J2-formalizacja} \\[3pt]

Emergencja U(1) z~symetrii fazy substratu
  & \statuslabel{Propozycja}
  & 5-krokowy dow\'od (thm:photon-emergence);
    warunkowo na aks.~\ref{ax:complex-substrate};
    \texttt{ex109} (12/12 PASS)
  & tw.~\ref{thm:photon-emergence}; \S\ref{sec:cechowanie} \\[3pt]

Emergencja SU(2) z~dw\'och ga\l{}\k{e}zi kiralnych
  & \statuslabel{Szkic}
  & mechanizm wskazany; brak pe\l{}nej derywacji
  & \S\ref{sec:cechowanie} \\[3pt]

Emergencja SU(3) z~fazy topologicznej
  & \statuslabel{Szkic}
  & mechanizm wskazany; brak pe\l{}nej derywacji
  & \S\ref{sec:cechowanie} \\[3pt]

Naruszenie CP: jedna faza Diracowska przy $N = 3$
  & \statuslabel{Propozycja}
  & $\pi_1(\mathbb{RP}^2)$; argument topologiczny
  & stw.~\ref{prop:CP-necessity} \\[3pt]

Forma ogona: $g(r){-}1 \sim (B\cos r + C\sin r)/r$
  & \statuslabel{Propozycja}
  & linearyzacja ODE przy $g=1$; sub.\ $w = r\cdot u$ (Dod.~K)
  & stw.~\ref{prop:K-tail-form}; dod.~\ref{app:wkb-atail} \\[3pt]

$A_{\rm tail} \propto (g_0{-}g^*)^4$: wyk\l{}adnik ze stopnia $V(g)$
  & \statuslabel{Propozycja}
  & $V(g)$ wielomian st.~4; $\alpha$-niezale\.zno\'s\'c (\texttt{ex62})
  & res.~\ref{res:K-SWKB}; dod.~\ref{app:wkb-atail} \\[3pt]

O-J1: analityczna formu\l{}a $A_{\rm tail}(g_0)$ (WKB~+~dopasowanie asympt.)
  & \statuslabel{Program}
  & $S_{\rm WKB} \propto (g_0{-}g^*)^{3/2}$; funkcja {\l}\k{a}cz\k{a}ca nieznana
  & prob.~\ref{prob:K-OJ1}; dod.~\ref{app:wkb-atail} \\[3pt]

O-J3: $\varphi$-FP + Koide $\to$ trzy masy leptonów ($r_{21}=206{,}77$,
  $r_{31}=3477{,}4$, $m_\tau = 1777{,}0$\,MeV, $+60$\,ppm od PDG)
  & \statuslabel{Propozycja [AN+NUM+POST($\mathbb{Z}_3$)]}
  & ODE substrat.\ + $\varphi$-FP $\to$ $r_{21}$;
    Koide $Q_K{=}3/2$ z~$\mathbb{Z}_3$ (za\l{}o\.zona, nie wyprowadzona);
    zero par.\ wolnych; hip.~,,harmoniczna'' $2g_0^e$ \textbf{obalona} ($6{,}9\%$)
  & prop.~\ref{prop:J-koide-chain}; \texttt{oj3\_tau\_selection.py} \\[3pt]

Precyzja $\varphi$-FP + Koide: $r_{31} = 3477{,}44$
  ($83$~ppm od PDG, \texttt{ex157})
  & \statuslabel{Propozycja [AN+NUM]}
  & rem.~\ref{rem:J-precision-ex157}; ODE substratowe $\textsf{rtol}=10^{-10}$
  & dod.~\ref{app:ogon-masy} \\[3pt]

$\varphi$-FP dzia\l{}a na kwarki: $r_{21}^{(d,s,b)} = 20{,}0$,
  $r_{21}^{(u,c,t)} = 588$ (zgodno\'s\'c z~PDG)
  & \statuslabel{Propozycja [NUM]}
  & res.~\ref{res:X-phiFP-universal}; Koide $K \neq 2/3$ dla kwark\'ow;
    $K$ jest ODE-niezmienniczy (wn.~\ref{cor:X-no-ODE-fix})
  & dod.~\ref{app:quark-sector}; \texttt{ex158}--\texttt{ex161} \\[3pt]

R11: $K(\varphi)=\varphi^2$ vs $\varphi^4$ --- rozwi\k{a}zany
  & \statuslabel{Twierdzenie}
  & linearyzacja $H_\Gamma$ ($\varphi^2$) vs pe\l{}na def.~$F_{\rm kin}$
    ($\varphi^4$); $H_\Gamma \neq F_{\rm kin}$; ODE substratowe ($\alpha{=}1$)
    preferowane (bariera duchowa w~$\alpha{=}2$ wyklucza $\varphi$-FP)
  & uw.~w~\S\ref{sec:ghost-resolution}; \texttt{ex163} \\[3pt]

$\alpha{=}1$ vs $\alpha{=}2$: rozstrzygni\k{e}cie
  & \statuslabel{Twierdzenie}
  & metryka $\Phi^{2/3}$ aksjomatyczna (A3); linearyzacja PPN
    niezale\.zna od~$\alpha$; $\kappa$, PPN, $N_e$, $n_s$ uniwersalne;
    ODE solitonu jedyny sektor zale\.zny --- $\alpha{=}1$ preferowane
  & wn.~\ref{cor:alpha1-preferred}; \texttt{ex166}, \texttt{ex167} \\[3pt]

Chiralna asymetria: $m_R/m_L = 2{,}03$ z~$V(g)$ i~$K_{\rm sub}$
  & \statuslabel{Propozycja}
  & $V'''(1) = {-}4$; $g_0^L(e) < g^*$ (bariera duchowa);
    \texttt{ex108} (9/9 PASS)
  & D.1d; OP-D1d-1 ZAMKNI\k{E}TY \\[3pt]

$n_s = 0{,}9662$ (MS numeryczny, $N_e = 60$)
  & \statuslabel{Twierdzenie}
  & Mukhanov-Sasaki z~$\nu = 3/2 + \varepsilon_H + \eta_H + 2\varepsilon_\psi$;
    \texttt{ex107} (10/10 PASS); 0{,}32$\sigma$ Planck
  & dod.~\ref{app:J2-formalizacja}; prop:starobinsky-emergence \\[3pt]

$r = 0{,}0033$; $r \cdot N_e^2 = 12$ (klasa Starobinsky'ego)
  & \statuslabel{Twierdzenie}
  & $r = 16\varepsilon_H$; BICEP/Keck: $r < 0{,}036$;
    \texttt{ex107}
  & prop:r-tensor \\

\bottomrule
\end{tabular}
\end{center}

% -------------------------------------------------------
\subsection*{Sektor kwantowy i~UV}
% -------------------------------------------------------

\begin{center}\footnotesize
\begin{tabular}{@{}p{5.2cm}lp{4.4cm}p{2.6cm}@{}}
\toprule
\textbf{Zdanie} & \textbf{Status} & \textbf{Warunki / za\l{}o\.zenia}
  & \textbf{Odsylacz} \\
\midrule

Substratowe obci\k{e}cie UV: $k_{\max} \sim 1/\lP$
  & \statuslabel{Twierdzenie}
  & wymiarowanie substratu ($\lP = $ min. skala)
  & stw.~\ref{prop:UV-cutoff} \\[3pt]

Jedno\-elementowa p\k{e}tla sko\'nczona:
$\delta m^2_{\rm sp}/m^2_{\rm sp} \sim \ell_P/({\ell_Y\Phi_0}) \sim 10^{-60}$
  & \statuslabel{Propozycja}
  & 1-p\k{e}tla z~$k_{\max}$; brak wy\.zszych rz\k{e}d\'ow
  & stw.~\ref{prop:one-loop-UV} \\[3pt]

Punkt sta\l{}y UV $(\tilde m^2_*, \tilde g_*) = (0, -2\pi^2/5)$
  & \statuslabel{Hipoteza}
  & beta-funkcje 1-p\k{e}tlowe; wymaga pełnego ERG
  & uw.~\ref{rem:IIIb-status} \\[3pt]

Asymptotyczne bezpiecze\'nstwo TGP
  & \statuslabel{Program}
  & wymaga pe\l{}nego r\'ownania Wetterika dla $K(\psi) = \psi^4$
  & uw.~\ref{rem:IIIb-status} \\

\bottomrule
\end{tabular}
\end{center}

% -------------------------------------------------------
\subsection*{Podsumowanie: iloraz predykcyjno\'sci}
% -------------------------------------------------------

\begin{center}\footnotesize
\begin{tabular}{@{}lrrl@{}}
\toprule
\textbf{Kategoria} & \textbf{Liczba zdań} & \textbf{\% pewności}
  & \textbf{Uwagi} \\
\midrule
Aksjomat              & 2  & --- & wej\'scie, nie dowodzone \\
Twierdzenie           & 23 & wysoka & +9 v41: metryka, $\kappa$, ghost, $\pi_1$, Fermi, $\varphi$-FP, $r_{21}$, $n_s$, $r$ \\
Propozycja            & 23 & \'srednia & +6: antypod., $S_{\rm TGP}$, profile, chiralnosc, U(1), $k{=}4$ wymiar., \l{}a\'ncuch entrop. \\
Hipoteza              & 5  & niska & $m_\tau$; T1+T2; $a_\Gamma\Phi_0{=}1$ (DESI DR2); CV${=}1$ \\
Szkic                 & 2  & --- & SU(2), SU(3) \\
Program               & 3  & --- & O-J1 (ogon), asympt.\ bezp., $\alpha_K{=}a_c^{-1}(a_\Gamma)$ \\
\midrule
\textbf{\'Lacznie}   & \textbf{58} & & +12 sesja v41; \'Sc.~9, $n_s$, $r$, chiralnosc; +2 v46: $k{=}4$, \l{}a\'ncuch \\
\midrule
Parametry Warstwy~II  & $N = 2$ & & $\Phi_0$, $a_\Gamma$
    ($\psi_{\rm ini}=7/6$: Warstwa~I$\to$III,
    stw.~\ref{prop:psi-ini-derived}, \texttt{ex86}) \\
Obserwacje wyja\'sniane & $M \geq 12$ & & patrz tab.~\ref{tab:param-classification} \\
Stosunek $M/N$        & $\geq 6$ & & \textbf{v35}: $M/N \geq 6$ (bez T1+T2);
    $M/N \geq 7$ (hip.~\ref{hyp:agamma-phi0-combo} $\Rightarrow N{=}1$) \\
\bottomrule
\end{tabular}
\end{center}

\begin{remark}[Priorytety redukcji niepewno\'sci --- stan po sesji v33+]\label{rem:status-priority}
\textbf{Sesja v33 zamkn\k{e}\l{}a}:
\begin{itemize}[nosep]
  \item Zadanie~B.2: forma metryki
    $g_{\mu\nu} = f(\Phi)\eta_{\mu\nu}$ awansowana
    z~\statuslabel{Hipoteza} do~\statuslabel{Twierdzenie} (forma)
    i~\statuslabel{Propozycja} (funkcja $f$);
    podstaw\k{a} jest tw.~\ref{thm:metric-form-uniqueness}
    + stw.~\ref{prop:g00-from-axioms}.
  \item Zadanie~N0-6: $m_{\rm sp} = \sqrt{\gamma}$ formalnie wyprowadzone
    z~N0-5 (stw.~\ref{prop:N0-6-from-N0-5}).
  \item Zadanie~C.2: tilt CMB awansowany do~\statuslabel{Propozycja}
    po numerycznej weryfikacji skryptem~\texttt{ex54}.
  \item \'Scie\.zka~8 ($r_{21}$ przez ZPE): \textbf{zamkni\k{e}ta} numerycznie
    (\texttt{ex53}); wymagane $x_A = -0{,}97 < 0$.
\end{itemize}
\textbf{Prze\l{}om sesji v33+} (skrypty \texttt{ex55}--\texttt{ex57}):
\begin{itemize}[nosep]
  \item \texttt{ex57} (skan szerokopasmowy z~elastycznym odbiciem przy~$g^*$):
    \[
      \left(\frac{A_{\rm tail}(g_0^\mu)}{A_{\rm tail}(g_0^e)}\right)^4
      = \left(\frac{A(2{,}00)}{A(1{,}24)}\right)^4 = 210{,}1
      \approx 206{,}77 \quad (\text{odch. }1{,}6\%)
    \]
    Globalne skalowanie: $M_{\rm raw} \sim g_0^{4{,}12}$ (potwierdzenie hipotezy $g^4$).
  \item Mechanizm: amplituda ogona oscylacyjnego $A_{\rm tail} = A\sin(r)/r$
    pe\l{}ni rol\k{e} sta\l{}ej sprz\k{e}\.zenia cz\k{a}stki z~substratem;
    $M \propto A_{\rm tail}^4$ t\l{}umaczy stosunek $r_{21}$.
\end{itemize}
\textbf{Wyniki \texttt{ex58} (warunek kwantowania $g_0$, sesja v33+)}:
\begin{itemize}[nosep]
  \item Hipoteza H1 ($B_{\rm tail}=0$, ,,czysty sinus''):
    pierwsze zero przy $g_0 = 1{,}2301 \approx g_0^e = 1{,}24$
    (odch.~$0{,}8\%$). \textbf{H1 potwierdzona dla elektronu.}
  \item Drugie zero $B_{\rm tail}$ przy $g_0 = 2{,}0797 \neq 2{,}00$;
    $(A(2{,}08)/A(1{,}23))^4 = 375 \neq 207$ ---
    sam warunek $B=0$ nie reprodukuje $r_{21}$.
  \item Obserwacja \textbf{z\l{}otej proporcji}:
    $g_0^\mu/g_0^e = 2{,}00/1{,}24 \approx \varphi = 1{,}618$;
    hipoteza: $g_0^e = \sqrt{5}-1$, $g_0^\mu = 2$
    ($g_0^\mu/g_0^e = \varphi$ dok\l{}adnie).
\end{itemize}
\textbf{Wyniki \texttt{ex59} (sektory topologiczne, sesja v33+)}:
\begin{itemize}[nosep]
  \item $M(g_0)$ ro\'snie monotonicznie: brak ekstremum przy $g_0 = 2{,}00$
    --- $g_0^\mu$ nie jest wyro\.znione przez $dM/dg_0 = 0$.
  \item Ka\.zdy sektor $n_{\rm bounce} = 0,1,2,\ldots$ ma \textbf{dok\l{}adnie jedno}
    zero $B_{\rm tail}$: $g_0 = 1{,}230\;(n=0)$, $2{,}080\;(n=1)$, $2{,}788\;(n=2)$.
  \item Wyk\l{}adniki skalowania: $n=0$: $\gamma=9{,}3$;
    $n=1$: $\gamma=4{,}4$ ($\approx 4$ !); $n=2$: $\gamma=3{,}4$.
  \item Najlepsza $g_0$ dla $r_{21}$: $g_0 = 1{,}982$,
    $(A/A_e)^4 = 205{,}4 \approx 207$ (odch.~$0{,}7\%$).
\end{itemize}
\textbf{Wyniki \texttt{ex60} (konwergencja $A_{\rm tail}$, sesja v33+)}:
\begin{itemize}[nosep]
  \item $A_{\rm tail}$ zbiega: zmienno\'s\'c $0{,}15\%$ (elektron), $0{,}63\%$ (mion)
    --- amplituda jest dobrze okre\'slona fizycznie.
  \item $g_0^\mu(r_{21}=207) = 1{,}989\pm0{,}005$, tj.~o~$1{,}1\%$ poni\.zej $g_0 = 2{,}00$.
  \item \textbf{Warunek z\l{}otej proporcji} (precyzja $0{,}04\%$):
    \[
      g_0^\mu(r_{21}=207)
      = \varphi \cdot z_0
      = 1{,}6180 \times 1{,}2301
      = 1{,}9903,
    \]
    gdzie $z_0 = 1{,}2301$ jest pierwszym zerem $B_{\rm tail}$ (elektron, sektor $n=0$).
  \item Interpretacja: kwantowanie TGP --- $g_0^e = z_0$ (H1), $g_0^\mu = \varphi\cdot z_0$
    --- reprodukuje $r_{21}\approx 207$ z dok\l{}adno\'sci\k{a} $< 0{,}1\%$
    (do analitycznego uzasadnienia).
\end{itemize}
\textbf{Wyniki \texttt{ex61} (weryfikacja warunku $\varphi$, sesja v33+)}:
\begin{itemize}[nosep]
  \item \textbf{Krytyczna diagnoza}: czyste $\varphi$-kryterium z~$g_0^e = z_0$:
    \[
      \left(\frac{A_{\rm inf}(\varphi z_0)}{A_{\rm inf}(z_0)}\right)^4
      = 252{,}07 \neq 207 \quad (\text{odch.~}21{,}9\%!)
    \]
    Interpretacja ex60 ($g_0^e = z_0$) jest \textbf{nieprecyzyjna}.
    Fizyczny elektron le\.zy przy $g_0^e = 1{,}24 \neq z_0 = 1{,}2301$.
  \item Mieszane kryterium: $(A(\varphi z_0)/A(1{,}24))^4 = 209{,}6$ (odch.~$1{,}35\%$)~---
    nadal poprawne.
  \item Seria generacji $\varphi^n z_0$: mion $n=1 \to 252$~(wz.~207),
    tauon $n=2 \to 5345$~(wz.~3477, odch.~$54\%$).
  \item Zale\.zno\'s\'c od $\alpha$: przy $\alpha = 2{,}5$ czyste $\varphi$ daje ratio~$= 197$
    ($5\%$ od~207); istnieje $\alpha^* \approx 2{,}3\text{--}2{,}4$ taki, \.ze ratio~$= 207$.
  \item Wra\.zliwo\'s\'c na $G_{\rm bounce}$: $z_0$ stale $= 1{,}2294$ --- efekt fizyczny,
    nie artefakt regularyzacji.
  \item Faza B--S: $\Delta\phi = 2{,}848\,\text{rad} = 0{,}907\pi$
    ($9{,}35\%$ od $\pi$).
\end{itemize}
\textbf{Wyniki \texttt{ex62} ($\varphi$-self-consistent fixed point, sesja v33+)}:
\begin{itemize}[nosep]
  \item \textbf{PRZE\L{}OM}: istnieje $g_0^* = 1{,}24771$ takie, \.ze
    \[
      \left(\frac{A_{\rm inf}(\varphi g_0^*)}{A_{\rm inf}(g_0^*)}\right)^4
      = 206{,}768 \quad (\text{odch.~}0{,}000001\%!)
    \]
    (warunek kwantowania \emph{self-consistent}, nie~H1 $B_{\rm tail}=0$).
  \item $\varphi g_0^* = 2{,}01884$ (kandydat $g_0^\mu$);
    $g_0^* - z_0 = +0{,}01761$;
    $g_0^* - 1{,}24 = +0{,}00771$.
  \item Brak prostego wzoru zamkni\k{e}tego dla~$g_0^*$
    (sprawdzono $\sqrt{5}-1$, $4/\pi$, $e^{-1/4}+0{,}47$, \ldots).
  \item $\alpha^* \approx 2{,}443$: przy tym $\alpha$ czyste $\varphi$ z~$z_0(\alpha^*)$ daje
    $r_{21}=207$ (TGP u\.zywa $\alpha=2$, odch. $= 0{,}44$).
  \item Tauon przy $g_0^*$: $(A(\varphi^2 g_0^*)/A(g_0^*))^4 = 4151$
    (eksp.~$3477$, odch.~$19\%$). Tauon FP: $g_0^{**}=1{,}2621$.
\end{itemize}
\textbf{Wyniki \texttt{ex63--ex66} (energetyka i~faza, sesja v33+)}:
\begin{itemize}[nosep]
  \item \textbf{ex63} (faza scatteringowa):
    $\delta_{\rm scatt}(g_0^*) = 0{,}015\,\text{rad} \approx 0$,
    $\delta(\varphi g_0^*) = 0{,}938\pi$; warunek B--S = warunek H1;
    $\varphi$-FP jest r\'ownaniem transcendentnym --- brak prostego kryterium fazowego.
  \item \textbf{ex64}: $M_{\rm raw} \propto A_{\rm tail}^{2{,}013}$ (nie $A^4$!);
    $M_{\rm raw}$ zdominowany przez ogon ($\sim A^2 R_{\rm max}$, rozbie\.zny).
    $\alpha^* = 2{,}441$, $|\alpha^* - \sqrt{6}| = 0{,}008$ ($0{,}3\%$).
  \item \textbf{ex65}: $M_{\rm kin} \propto A^{2{,}013}$, $|M_{\rm pot}| \propto A^{2{,}000}$,
    $M_{\rm pot} < 0$; $M_{\rm tot} = M_{\rm kin}+M_{\rm pot}$ zmienia znak w~sektorze $n=0$
    (elektron: $-0{,}22$, mion: $+7{,}44$). Ogon: $M_{\rm tail} \approx 0$ (kasowanie).
  \item \textbf{ex66} + \textbf{wniosek analityczny}:
    Tryb $k=1$ jest \textbf{trybem zerowym} zlinearyzowanego ODE TGP:
    \[
      \delta'' + \frac{2}{r}\delta' + \delta = 0
      \implies \mathcal{E}_{\rm linear} = \frac{A^2\cos(2r)}{2r^2} \to 0
    \]
    Energia trybu zerowego $= 0$ dok\l{}adnie (kasowanie kin.+pot. w~ogonie).
    Fizyczna masa $= O(A^4)$ nieliniowe korekcje:
    \[
      M_{\rm TGP} \approx c\,A_{\rm tail}^4 + O(A_{\rm tail}^6)
    \]
    \textbf{To jest matematyczne uzasadnienie $M \propto A_{\rm tail}^4$
    z~aksjomatu TGP} --- status: \statuslabel{Propozycja}.
\end{itemize}
\textbf{Wyniki \texttt{ex69} (selekcja $g_0^*$, sesja v33++), 4/5 test\'ow}:
\begin{itemize}[nosep]
  \item $\delta(\alpha) \equiv g_0^*(\alpha) - z_0(\alpha)$
    jest \textbf{liniowe} w~$\alpha$ ($R^2 = 0{,}99993$):
    \[
      \delta(\alpha) = -0{,}044063\,\alpha + 0{,}107667,
      \qquad \delta = 0 \;\text{przy}\; \alpha = 2{,}4435.
    \]
    Przy $\alpha = \alpha^* = 2{,}441$ warunki H1 ($B_{\rm tail}=0$)
    i~$\varphi$-FP \textbf{zbiegaj\k{a} si\k{e}} w~jeden punkt:
    $g_0^*(\alpha^*) = z_0(\alpha^*) = 1{,}2585$ (r\'o\.znica $0{,}00011$).
  \item Implikacja: $\alpha = \alpha^* \Rightarrow r_{21} = f(z_0(\alpha^*)) = 207{,}05$
    \textbf{czysto z~warunku H1}, bez wej\'scia empirycznego!
  \item Hipoteza: $\alpha^* = \sqrt{6} \approx 2{,}449$ (0{,}3\% od 2{,}441)?
\end{itemize}
\textbf{Bie\.z\k{a}ce priorytety (P1)}:
\begin{enumerate}[nosep]
  \item Dlaczego $\alpha_{\rm TGP} = 2$ (nie $\alpha^* = 2{,}441$)?
    Je\'sli TGP wewn\k{e}trznie wymaga $\alpha = \alpha^*$, to $r_{21}$
    jest \textbf{przewidziany} z~aksjomatu H1 + $\varphi$-skalowanie.
  \item \textbf{ex71: dwa zera $f_{z_0}(\alpha)=r_{21}$}
    (\statuslabel{Propozycja}; \textbf{korekta po ex84}):
    R\'ownanie $f(z_0(\alpha),\alpha) = r_{21}$ ma dok\l{}adnie dwa
    rozwi\k{a}zania przy $\beta=1$:
    \[
      \alpha_1^* \approx 2{,}440, \qquad \alpha_2^* \approx 2{,}695.
    \]
    (\emph{Uwaga}: seria ex71--ex83 podawa\l{}a $\alpha_1^*=2{,}44142$,
    $\alpha_2^*=2{,}74175$ --- s\k{a} to warto\'sci \textbf{b\l{}\k{e}dne},
    wynikaj\k{a}ce z~artefaktu grubej siatki \texttt{ex71};
    poprawione przez \texttt{ex84}, patrz ni\.zej.)
    Przy obu $\alpha^*$: H1~($B_{\rm tail}=0$) $\equiv$ $\varphi$-FP,
    co daje $r_{21}=206{,}768$.
  \item \textbf{ex84: \textsc{Falsyfikacja} formu\l{}y $S=2\pi-\tfrac{11}{10}$}
    (\statuslabel{Twierdzenie} --- wynik negatywny):
    Bezpo\'srednia weryfikacja $f_{z_0}(\alpha) = r_{21}$
    bez hardcoded warto\'sci (\texttt{ex84}, 1/5~PASS):
    \begin{itemize}[nosep]
      \item Prawdziwe zera: $\alpha_1^* = 2{,}4396$, $\alpha_2^* = 2{,}6953$
        (NIE $2{,}4414$ i~$2{,}7418$ jak podawa\l{}a seria ex71--ex83).
      \item Suma $S = \alpha_1^*+\alpha_2^* = 5{,}1349$,
        odchylenie \textbf{9307\,ppm} od $2\pi-11/10 = 5{,}1832$.
      \item \'{Z}r\'od\l{}o artefaktu: \texttt{ex71} u\.zy\l{} 25-punktowej siatki
        $z_0(\alpha)$ --- spurious zero crossings przepropagowane
        przez ca\l\k{a} seri\k{e} ex72--ex83 jako hardcoded truth.
      \item Zachowane: dwa zera fz0 przy $\alpha^*\approx 2{,}440$ i~$2{,}695$
        s\k{a} realne (P1, P2~PASS); punkt sta\l{}y $\varphi$-FP przy
        $g_0^*=1{,}24771$ jest realy.
      \item Odrzucone: formu\l{}a $S=2\pi{-}11/10$,
        generalizacja $S(n)=2\pi{-}(5n+1)/10$,
        $D(n)=(n+1)/10$ --- wszystkie bazowa\l{}y na b\l{}\k{e}dnych $\alpha^*$.
    \end{itemize}
    \textbf{\'Scie\.zka~8--9 (ogon algebraiczny)}: \textsc{Zamkni\k{e}ta}.
  \item \textbf{ex73--ex77: $\beta=1$ jako centrum symetrii}
    (\statuslabel{Hipoteza}; \textbf{wymaga re-weryfikacji}):
    Przy starych $\alpha^*$ (b\l{}\k{e}dnych): dwu-zerowa struktura
    $\delta(\alpha)$ wyst\k{e}puje wy\l\k{a}cznie dla $\beta=1$;
    $\beta=1$ by\l{}o centrum okna bifurkacji z~dok\l{}adno\'sci\k{a}
    $1{,}87\times10^{-5}$.
    \textbf{Status po ex84}: wynik wymaga re-weryfikacji
    z~poprawnymi warto\'sciami $z_0(\alpha)$ (poprawna siatka, $>100$~pkt).
    Je\'sli $\beta=1$ pozostaje centrum symetrii przy poprawnych $\alpha^*$
    $\Rightarrow$ aksjomat N0-5 ($\beta=\gamma$) ma dodatkowe uzasadnienie
    dynamiczne.
  \item \textbf{ex85: Nowe relacje algebraiczne Warstwy~II}
    (\statuslabel{Hipoteza}; \texttt{ex85}, 5/6~PASS):
    Weryfikacja z~danymi Planck~2018, DESI~DR1+CMB:
    \begin{itemize}[nosep]
      \item \textbf{T1}: $a_\Gamma\cdot\Phi_0 = 1$ ---
        dok\l{}adno\'s\'c $0{,}09\%$ ($-0{,}12\sigma$) z~DESI+CMB.
      \item \textbf{T2}: $r_{21} = \Phi_0 \cdot \alpha_K$ ---
        dok\l{}adno\'s\'c $0{,}048\%$ z~Planck~2018.
      \item \textbf{Konsekwencja T1+T2}:
        $a_\Gamma = (\alpha_K \sqrt{r_{21}})^{-2/3} = 0{,}04011$
        (0{,}15\% od fizycznego $0{,}040049$).
        Je\'sli obie relacje dok\l{}adne:
        $a_\Gamma$ znika z~Warstwy~II
        $\Rightarrow$ $N_{\rm param}=2$ ($\Phi_0, \psi_{\rm ini}$).
    \end{itemize}
  \item \textbf{Sesja v34 (Dodatek~K: analityczna asymptotyka ogona i~predykcja tauonu)}:
    \begin{itemize}[nosep]
      \item \textbf{Prop.~\ref{prop:K-tail-form}} (stw.~K-tail-form):
        forma ogona $g(r){-}1 \sim (B\cos r + C\sin r)/r$
        formalnie dowiedziona przez linearyzacj\k{e} ODE przy $g=1$
        i~podstawienie $w = r\cdot u$; awans do~\statuslabel{Propozycja}.
      \item \textbf{$A_{\rm tail} \propto (g_0{-}g^*)^4$}
        (res.~\ref{res:K-SWKB}):
        wyk\l{}adnik $n \approx 4$ wynika ze stopnia potencja\l{}u $V(g)$
        (wielomian st.~4), niezale\.znie od~$\alpha$
        (potwierdzono numerycznie skryptem~\texttt{ex62} dla
        $\alpha \in [1{,}0;\, 3{,}0]$); awans do~\statuslabel{Propozycja}.
      \item \textbf{O-J1} (prob.~\ref{prob:K-OJ1}):
        analityczna formu\l{}a $A_{\rm tail}(g_0)$ przez metod\k{e} dopasowania
        asymptotycznego --- ca\l{}kowanie WKB daje
        $S_{\rm WKB} \propto (g_0{-}g^*)^{3/2}$, ale {\l}\k{a}czenie
        z~obszarem ogonowym jest nieznane;
        pozostaje \statuslabel{Program}.
      \item \textbf{O-J2} ($\varphi$-FP): {\bfseries ZAMKNI\k{E}TY}.
        $g_0^\mu = \varphi \cdot g_0^*$; $(A_\mu/A_e)^4 = 206{,}768$
        (dok\l{}adna zgodno\'s\'c z~PDG) potwierdzone we
        wszystkich trzech ODE (uproszcz., pe\l{}ne, substrat.);
        status~\statuslabel{Twierdzenie}.
      \item \textbf{O-J3} (rem.~\ref{rem:K-substrate-tau}):
        ODE substratowe ($K_{\rm sub} = g^2$, ghost-free):
        $g_0^\tau = 1{,}729$, $m_\tau = 1776{,}8$~MeV ($-14$~ppm),
        Koide $= 0{,}66666$ ($-11$~ppm);
        zasada selekcji $g_0^\tau$ otwarta ($g_0^\tau/g_0^e \approx 2$
        z~odch.\ $1{,}1\%$); status~\statuslabel{Propozycja}.
      \item Komplet skrypt\'ow: \texttt{ex60} (topologia sektora ogona),
        \texttt{ex61} (predykcja tauonu, kandydaci C1--C5),
        \texttt{ex62\_phi\_fixed\_point.py} (skan kinetycznego~$\alpha_K$,
        wyznaczenie $\alpha_1^*, \alpha_2^*$, $\varphi$-FP)
        --- udokumentowane w~Dod.~\ref{app:wkb-atail}.
    \end{itemize}

  \item \textbf{Sesja v34 --- Dodatek~L (\'{S}cie\.zka~10: formu\l{}a sumy $\alpha_1^*+\alpha_2^*$)}:
    \begin{itemize}[nosep]
      \item \textbf{Definicja} (def.~L-F):
        $F(\alpha_{\rm kin}) = \bigl(A_{\rm tail}(\varphi z_0) / A_{\rm tail}(z_0)\bigr)^4$;
        dwa pierwiastki $F(\alpha^*) = r_{21}$:
        $\alpha_1^* \approx 2{,}440$, $\alpha_2^* \approx 2{,}695$
        (poprawione przez \texttt{ex84}).
      \item \textbf{Hipoteza sumy} (hyp.~\ref{hyp:L-sum-formula}):
        Pierwotna formu\l{}a $\alpha_1^* + \alpha_2^* = 2\pi - \tfrac{11}{10}$
        pochodzi\l{}a z~artefaktu siatki serii ex71--ex83
        ($\alpha_1^*=2{,}441$, $\alpha_2^*=2{,}742$ --- warto\'sci b\l{}\k{e}dne).
        Po korekcie (\texttt{ex84}): $S = 5{,}1349$, odch.\ $9307\,\text{ppm}$
        od $2\pi{-}11/10$ --- formu\l{}a \textbf{obalona}.
        Status: \statuslabel{Hipoteza} (obalona).
      \item \textbf{Zachowane wyniki pozytywne}:
        istnienie dw\'och zer $F(\alpha^*)=r_{21}$ (\statuslabel{Propozycja});
        punkt sta\l{}y $\varphi$-FP przy $g_0^* = 1{,}24771$
        (\statuslabel{Propozycja});
        interpretacja fizyczna $\alpha^*$ jako efektywnych sprz\k{e}\.ze\'n
        kinetycznych sektor\'ow leptonowych (\statuslabel{Hipoteza}).
      \item \textbf{O-L1 zamkni\k{e}te} (\texttt{ex88}, v36, 4/4~PASS):
        Okno $B_{\rm coeff}$ \textbf{sta\l{}e} $[28,42]$ (niezale\.zne od
        $r_{\rm max}$) --- amplituda niezale\.zna od $r_{\rm max}$ przy
        sta\l{}ym oknie fizycznym.
        Wyniki finalne:
        $\alpha^*_{1,\infty} = 2{,}4318$,
        $\alpha^*_{2,\infty} = 2{,}6356$ ($\sigma=0$, dok\l{}adno\'s\'c maszynowa).
        $S(\infty) = 5{,}0674 \neq 2\pi - 11/10 = 5{,}1832$ (odch.\ $22336\,\text{ppm}$)
        --- formu\l{}a sumy \textbf{definitywnie obalona}.
        Brak prostej formu\l{}y algebraicznej dla $S$.
        O-L2--O-L5 otwarte.
        Skrypty: \texttt{ex62}--\texttt{ex88}.
    \end{itemize}

  \item \textbf{Sesja v35 --- F5 zamkni\k{e}te:
    atraktor $\psi_{\rm ini} = 7/6$} (\statuslabel{Propozycja};
    \texttt{ex86}, 9/9~PASS):
    \begin{itemize}[nosep]
      \item \textbf{T1--T2} (algebraika): $W(7/6) = 0$,
        $W'(7/6) = -49\gamma/18 < 0$ --- stabilny atraktor slow-roll.
        Cz\k{e}sto\'s\'c: $\omega_{\rm cosmo}/H_0 \approx 8{,}19$
        (niedot\l{}umione oscylacje przy $z=0$).
      \item \textbf{T3--T4} (zamro\.zenie): $\eta_{\rm BBN} = 1{,}2\times10^{-31}$,
        $\Delta\psi_{\rm rad} \sim 3\times10^{-11}$ ---
        pole absolutnie zamro\.zone w~erze radiacji.
      \item \textbf{T5} (BBN): $\psi_{\rm ini} = 7/6
        \;\Rightarrow\; |\Delta G/G| = 0{,}000\%$ ---
        \textbf{jedyna warto\'s\'c} daj\k{a}ca doskonale
        spe\l{}nione ograniczenie BBN.
        Por.\ $\psi_{\rm ini} = 1{,}0 \;\Rightarrow\; |\Delta G/G| = 13{,}7\%$.
      \item \textbf{Awans:} \texttt{prop:psi-ini-derived}
        $\mathrm{Hipoteza} \to \mathrm{Propozycja}$;
        $N_{\rm param}:\ 3 \to 2$
        ($\psi_{\rm ini}$ opuszcza Warstw\k{e}~II).
      \item \textbf{Nowa predykcja} (v35): oscylacje $G_{\rm eff}(t)$
        z~$\omega_{\rm cosmo}/H_0 \approx 8{,}19$ dla
        $|\psi_{\rm ini} - 7/6| \neq 0$ --- testowalne przez
        projekty pomiaru $G(t)$ (LLR, GRACE-FO).
    \end{itemize}

  \item \textbf{Sesja v35 --- F6 zamkni\k{e}te:
    napi\k{e}cie $r \approx 0{,}27$ vs $r < 0{,}036$ ---
    artefakt b\l{}\k{e}dnego wzoru}:
    \begin{itemize}[nosep]
      \item \textbf{Diagnoza}: \texttt{rem:min-observables}
        u\.zywa\l{}o relacji kanonicznej $r = 16\varepsilon_\psi$
        z~$\varepsilon_\psi = (1-n_s)/2 \approx 0{,}017$,
        co~daje $r \approx 0{,}27$.
        Jest to parametr $\eta$ (dominuj\k{a}cy w~Starobinskim),
        \textbf{nie} $\varepsilon$ u\.zywane w~relacji $r = 16\varepsilon$.
      \item \textbf{Poprawka}: TGP nale\.zy do klasy Starobinsky'ego
        (stw.~\ref{prop:r-tensor}, dod.~G);
        $\varepsilon = 3/(4N_e^2)$,
        $r = 12/N_e^2 \approx 0{,}003$ ($N_e=60$) ---
        spe\l{}nia $r < 0{,}036$ \textbf{z~du\.zym marginesem}.
      \item \textbf{$\sigma_{ab}$ a inflacja}: $\sigma_{ab} = 0$
        podczas inflacji (substrat w~fazie symetrycznej;
        $K_{ab} \propto \delta_{ab}$, def.~\ref{def:sigma-ab}).
        Pole $\sigma_{ab}$ aktywne wy\l{}\k{a}cznie
        po~podgrzaniu --- dla \'zr\'ode\l{} anizotropowych.
      \item \textbf{Predykcja}: $r \approx 0{,}003$
        (odr\'o\.znialne od Higgsinflacji/Starobinskiego
        przez CMB-S4 dopiero przy $r \sim 0{,}001$;
        LiteBIRD: granica $r < 0{,}002$).
    \end{itemize}

  \item \textbf{Sesja v36 --- O-L1b zamkni\k{e}te:
    finalne $\alpha^*_\infty$ bez zale\.zno\'sci od $r_{\rm max}$}
    (\texttt{ex88}, 4/4~PASS):
    \begin{itemize}[nosep]
      \item \textbf{G\l{}\'owne odkrycie}: stałe okno $B_{\rm coeff}\in[28,42]$
        eliminuje ca\l{}kowicie zale\.zno\'s\'c wynik\'ow od $r_{\rm max}$
        ($r_{\rm max} \in \{100,120,150,200,300\}$ --- identyczne $\alpha^*$
        z~dok\l{}adno\'sci\k{a} maszynow\k{a}).
        Artefakt \texttt{ex87v3} potwierdzony jako skutek skalowanego okna.
      \item \textbf{Finalne warto\'sci}:
        $\alpha^*_{1,\infty} = 2{,}43183767$,
        $\alpha^*_{2,\infty} = 2{,}63557742$,
        $S(\infty) = 5{,}06741509$ (cf.\ $2\pi - 11/10 = 5{,}18319$,
        odch.\ $22336\,\text{ppm}$).
      \item \textbf{Obalenie formu\l{}y sumy:} \textbf{ostateczne} ---
        nie tylko przy $r_{\rm max}=120$ ($9307\,\text{ppm}$, \texttt{ex84}),
        ale dla ca\l{}ego zakresu $r_{\rm max}$; warto\'s\'c $S=5{,}0674$
        stabilna i~bez prostej reprezentacji algebraicznej.
      \item \textbf{F7 zamkni\k{e}te analitycznie}:
        $n_s = 1 - 2/N_e$ (Starobinsky, rem:F7-ns-precision);
        TGP zdegenerowane z~Starobinskim przy precyzji Planck~2018;
        korekcja $\varepsilon_\psi = O(N_e^{-2}) \approx 0{,}0002$ poni\.zej
        progu CMB-S4.
      \item \textbf{$\sigma_{ab}=0$ podczas inflacji} (rem:sigma-inflation):
        substrat w~fazie symetrycznej $\to K_{ab}\propto\delta_{ab}$
        $\to \sigma_{ab}=0$ --- predykcja testowalna LiteBIRD
        $r\approx 0{,}003\pm 0{,}001$ (\statuslabel{Propozycja}).
      \item \textbf{O-L2 zamkni\k{e}te} (\texttt{ex89}, v36):
        skan $\beta_{\rm pot}\in\{0{,}5;\,0{,}8;\,1{,}0;\,1{,}2;\,1{,}5\}$
        --- dwa zera $F(\alpha^*)=R_{21}$ istniej\k{a}
        \textbf{wy\l{}\k{a}cznie dla $\beta_{\rm pot}=1$}
        (fizyczne N0-5); zera nie s\k{a} topologiczne.
        hyp:L-alpha-eff: bez awansu.
      \item \textbf{O-L3 zamkni\k{e}te cz\k{e}\'sciowo} (\texttt{ex91}, v36,
        4/5~PASS): \'Scie\.zka~9 ($\alpha_{\rm TGP}=2$) daje
        $F(2)=263{,}9\neq R_{21}$ --- unifikacja \textbf{cz\k{e}\'sciowa};
        \'scie\.zki 9 i~10 s\k{a} komplementarne, nie identyczne.
    \end{itemize}

  \item \textbf{Sesja v37 --- O-L4 i~O-L5: profil $F(\alpha)$ i~sektor tau}
    (\texttt{ex92}, 2026-03-28):
    \begin{itemize}[nosep]
      \item \textbf{Dwie ga\l{}\k{e}zie $z_0(\alpha)$}:
        ga\l{}\k{a}\'z g\'orna ($z_0 \approx 2{,}3$--$2{,}5$, $\alpha < 1{,}4$,
        niefizyczna) i~ga\l{}\k{a}\'z fizyczna ($z_0\approx 1{,}11$--$1{,}26$,
        $\alpha > 1{,}5$; u\.zywana w~\texttt{ex88}).
      \item \textbf{O-L4 --- mechanizm $\alpha^*\neq\alpha_{\rm TGP}$}:
        Na ga\l{}\k{e}zi fizycznej $F(\alpha)$ maleje od $F(2)=264{,}1$
        do minimum $F_{\min}\approx 201{,}0$ przy $\alpha_{\min}\approx 2{,}54$,
        po czym ro\'snie. $\alpha_{\rm TGP}=2$ le\.zy po~\textbf{lewej}
        stronie minimum (ga\l{}\k{a}\'z malej\k{a}ca).
        Dwa zera: $\alpha^*_1\approx 2{,}43$ (F maleje przez $R_{21}$)
        i~$\alpha^*_2\approx 2{,}64$ (F ro\'snie przez $R_{21}$),
        spójne z~\texttt{ex88} do $\pm0{,}002$.
        Otwarte: analityczna formu\l{}a $\alpha_{\min}$.
      \item \textbf{O-L5 --- $\alpha^*_3$ potwierdzona} (\texttt{ex93}+\texttt{ex94}, v37):
        skan $\alpha\in[2{,}70;\,3{,}30]$, krok $0{,}01$, Pool(6), czas $2333\,\text{s}$.
        $G_3(\alpha)$ monotonically malejąca (korekta \texttt{ex92}).
        Jedno zero $G_3=R_{31}$:
        \textbf{$\alpha^*_3 = 2{,}929524$},
        $|\Delta\alpha^*_3|=2{,}3\times10^{-8}$ dla $r_{\rm max}\in\{120\ldots300\}$.
        $G_3(\alpha^*_3)=3495{,}2$ ($+0{,}52\%$ od $R_{31}$; szum $\pm20$ j.).
        \textbf{Formu\l{}a sumy}:
        $S_3=\alpha^*_1+\alpha^*_2+\alpha^*_3=7{,}9969\approx8=4\alpha_{\rm TGP}$
        ($-383\,\text{ppm}$; kandydat, w granicach szumu $G_3$).
        $\alpha^*_{\tau\mu}=2{,}8534$: $G_3/F=R_{32}$, $\sigma\approx0$.
      \item Skrypt \texttt{ex92}: $N_\alpha=120$, $r_{\rm max}=120$,
        \texttt{Pool}(6), czas $484\,\text{s}$; 116/120 punkt\'ow poprawnych.
    \end{itemize}

  \item \textbf{Sesja v38 --- \textsc{korekta O-L5}: artefakt okien ekstrapolacji}
    (\texttt{ex95}, 2026-03-28):
    \begin{itemize}[nosep]
      \item \textbf{Krytyczny b\l{}\k{a}d wykryty}: $\varphi^2 z_0\approx3{,}31$
        --- soliton taonowy startuje z~du\.zy $g_0$;
        strefa bliska ($r<30$) nie jest asymptotyczna.
        Okno $[20,36]$ daje $A_{\rm tail}=2{,}838$ zamiast prawdziwego
        $\approx2{,}994$ (okna dalekie $[60,136]$).
      \item \textbf{Skutek}: standardowe okna $[20,66]$
        zani\.zaj\k{a} $A_{\rm tail}(\varphi^2 z_0)$
        $\Rightarrow$ $G_{3,\rm std}\approx3410$ (fa\l{}szywnie niskie)
        vs.\ $G_3^{\rm true}\approx4860$--$5000$.
        Prawdziwe $G_3 > R_{31}=3477$ dla wszystkich
        $\alpha\in[2{,}90;\,2{,}95]$.
      \item \textbf{O-L5 --- korekta}: $\alpha^*_3=2{,}929524$ i~$S_3=8$
        z~\texttt{ex93}/\texttt{ex94} s\k{a} \textbf{artefaktami}.
        Prawdziwe $\alpha^*_3$ nieznane.
        Zaplanowany \texttt{ex96}: skan $G_3(\alpha)$ z~oknami
        $[60,136]$ na $\alpha\in[2{,}5;\,5{,}0]$.
      \item Skrypt \texttt{ex95}: $r_{\rm max}=150$, 3 zestawy okien
        (std/ext/far), \texttt{Pool}(6), czas $893\,\text{s}$.
    \end{itemize}

  \item \textbf{Sesja v38 c.d. --- O-L5: skan szeroki i test zbieżności}
    (\texttt{ex96}+\texttt{ex97}, 2026-03-28):
    \begin{itemize}[nosep]
      \item \textbf{\texttt{ex96} --- skan $G_3(\alpha)$ z oknami FAR $[60,136]$
        na $\alpha\in[2{,}5;\,5{,}0]$}:
        $G_{3,\rm far} > R_{31}=3477$ \textbf{wszędzie}.
        Minimum $G_{3,\rm far}=4667$ przy $\alpha=2{,}50$ ($+34\%$ nad $R_{31}$).
        Brak zmiany znaku $G_3-R_{31}$ na całym zakresie.
        Dodatkowe: $\alpha^*_{2,\rm far}=2{,}7538$ ($F_{\rm far}=R_{21}$;
        vs.\ $\alpha^*_2=2{,}6356$ ze~STD --- przesunięcie $+0{,}118$).
      \item \textbf{\texttt{ex97} --- test zbieżności $A_{\rm tail}$ do $r_{\rm max}=300$}:
        Per-okno $A_\tau\approx2{,}994$--$2{,}998$ stabilne od $r=60$ do $r=286$.
        Trend $dA/dr=+9\times10^{-6}/r$ --- plateau potwierdzone.
        $G_3^{\rm true}(\alpha^*_3)\approx5558$ ($+60\%$ nad $R_{31}$).
        \textbf{Wniosek definitywny}: formuła $G_3=(A_\tau(\varphi^2 z_0)/A_e(z_0))^4=R_{31}$
        \textbf{nie ma rozwiązania} na gałęzi fizycznej.
        Sektor $\tau$ wymaga innej parametryzacji.
        O-L5: \emph{otwarte ponownie}; hipoteza $S_3=8$ bez podstawy numerycznej.
    \end{itemize}

  \item \textbf{Sesja v39 --- O-L5: systematyczna falsyfikacja scenariuszy $\tau$}
    (\texttt{ex100}--\texttt{ex102}, 2026-03-29):
    \begin{itemize}[nosep]
      \item \textbf{\texttt{ex100} --- S1: $G_{3,\rm far}$ na $\alpha\in[1{,}5;\,2{,}55]$}:
        Minimum $G_{3,\rm far}=4511$ przy $\alpha=2{,}35$ ($+30\%$ nad $R_{31}$).
        Brak zmiany znaku. $G_3(\varphi^3 z_0) \gg R_{31}$ wszędzie.
        Precyzyjne $\alpha_{\min,\rm far}=2{,}5595$ (brentq $dF/d\alpha=0$);
        $F_{\min,\rm far}=196{,}85 < R_{21}=206{,}77$.
        \textbf{S1 = falsyfikowany}.
      \item \textbf{\texttt{ex101} --- S2: rezonans czy fizyka przy $\alpha\approx7{,}1$?}
        Per-okno profil $A_\tau$: okno $[60,76]$ daje $A_\tau\approx3{,}09$
        przy $\alpha=7{,}0$, podczas gdy okna $[80,136]$ dają
        $A_\tau\approx4{,}35$--$4{,}67$ (plateau). Prawdziwe
        $G_3(\alpha=7{,}0)\approx65000 \gg R_{31}$.
        Zmiany znaku przy $\alpha\approx6{,}77$ i~$\alpha\approx7{,}36$
        to \textbf{artefakty biasu okna $[60,76]$}.
        Stabilno\'s\'c $R_{\rm max}$ (150 vs.\ 300) nie wyklucza artefaktu okiennego.
        \textbf{S2 = falsyfikowany}.
      \item \textbf{\texttt{ex102} --- S4+S5: inne hierarchie mno\.znikowe}
        (okna VFAR $[80,156]$, 11~mnożników):
        Najbli\.zej $R_{31}$: $m=2{,}5$ przy $\alpha=\alpha^*_{1,\rm far}$
        daje $G_3=3519{,}2$ ($+42$; $+1{,}2\%$) --- brak zera.
        Skan $G_{3,\rm vfar}(\varphi^2 z_0)$ na $\alpha\in[2{,}3;\,5{,}0]$:
        minimum $G_{3,\rm vfar}=4519$ przy $\alpha=2{,}30$ ($+30\%$),
        brak zmiany znaku.
        $\alpha_{\min,\rm vfar}=2{,}5574$ (VFAR brentq; vs.~FAR:~$2{,}5595$, $\Delta=0{,}002$).
        $F_{\rm far}$ vs.\ $F_{\rm vfar}$: $|\Delta F|\le0{,}28$ --- okna nieistotne dla mionu.
        \textbf{S4 = falsyfikowany; S5 = wykonane}.
      \item \textbf{Wniosek globalny v39}:
        Formu\l{}a $G_3=(A_\tau(m\cdot z_0)/A_e(z_0))^4=R_{31}$ nie ma rozwi\k{a}zania
        dla \.zadnego z~testowanych mno\.znik\'ow
        $m\in\{\varphi,\sqrt{2},\sqrt{3},\sqrt{5},e,2,2{,}5,\varphi^2,\varphi^3,3,1{+}\varphi\}$
        na ga\l{}\k{e}zi fizycznej $\alpha\in[2{,}3;\,5{,}0]$.
        Sektor $\tau$ nie jest opisany formu\l{}\k{a} amplitudow\k{a} TGP.
        \textbf{O-L5: definitywnie otwarte --- wymaga nowej parametryzacji}.
    \end{itemize}

  \item \emph{Hipoteza OP-3}: $\alpha_K = a_c^{-1}(a_\Gamma)$.
  \item \emph{Sektory cechowania} SU(2), SU(3).
\end{enumerate}
\end{remark}

% ================================================================
% Aktualizacja v47 (2026-04-09) — nowe wyniki
% ================================================================
\begin{remark}[Aktualizacja v47 --- nowe domkni\k{e}cia
  i~korekty sp\'ojno\'sci]\label{rem:status-v47}
\begin{enumerate}[nosep]
  \item \textbf{Element obj\k{e}to\'sciowy:}
    Wszystkie odwo\l{}ania do starego
    $\sqrt{-g_{\mathrm{eff}}} = \psi^4$ w~sek08\_formalizm.tex
    poprawione na $c_0\psi$
    (5 poprawek, zob.\ sek.~\ref{sec:akcja-zunifikowana}).
    R\'ownanie FRW (prop.~\ref{prop:FRW-derivation})
    zaktualizowane: $2\dot\psi^2/\psi \to 3\dot\psi^2/\psi$.

  \item \textbf{M1 --- jednoznaczno\'s\'c $k = 4$ (stw.~\ref{thm:R-k4-definite}):}
    Nowe twierdzenie w~dodatku~R:
    fizyczna masa solitonu w~$d = 3$ ma jednoznaczne
    rozwini\k{e}cie $M = c_4 A_{\rm tail}^4 + \mathcal{O}(A^6)$,
    poniewa\.z rz\k{a}d $n = 3$ jest logarytmicznie rozbie\.zny
    (schematycznie zale\.zny) --- unikalne dla $d = 3$.
    Status: \statuslabel{Twierdzenie}~[AN+NUM].

  \item \textbf{M2 --- $\lambda_{\mathrm{eff}} = 2/\PhiZero^4$
    (prop.~\ref{prop:I-lambda-Z2}):}
    Nowa propozycja w~dodatku~I: czynnik $\mathbb{Z}_2$
    zamyka luk\k{e} mi\k{e}dzy estymacj\k{a} wymiarow\k{a}
    a~skanem solitonowym ($1{,}7\%$ od $\lambda^*$).
    Status: \statuslabel{Propozycja}~[AN+NUM].

  \item \textbf{M3 --- $Q_K = 3/2$ z~$N_{\rm gen} = 3$
    (stw.~\ref{thm:T-QK-from-Ngen}):}
    Nowe twierdzenie w~dodatku~T:
    $Q_K = 2N/(N+1) = 3/2$ wynika z~$N = 3$ generacji.
    \textbf{Awans w~v47b}: hipoteza ekwipartycji
    $r = \sqrt{N{-}1}$ okazuje si\k{e} \textbf{algebraiczn\k{a}
    konsekwencj\k{a}} $Q_K = 3/2$ w~parametryzacji Brannena
    ($r^2 = 2(N/Q_K - 1) = 2$,
    prop.~\ref{prop:T-equipartition-info}).
    Status: \statuslabel{Twierdzenie}~[AN+NUM],
    T-OP1 \textbf{zamkni\k{e}ty}.

  \item \textbf{M4 --- sektor kwarkowy
    (prop.~\ref{prop:X-m0-confinement}--\ref{prop:X-universal-A},
    tw.~\ref{thm:X-A-golden}):}
    Nowe propozycje w~dodatku~X:
    masa konfinementowa $m_0$ z~potencja\l{}u liniowego;
    przesuni\k{e}ty Koide $K(m_k + m_0) = 3/2$;
    uniwersalna sta\l{}a $\mathcal{A} = a_\Gamma/\varphi = 0.02472$
    (odch.\ $0.31\%$). R12 istotnie zaw\k{e}\.zony.
    Status: \statuslabel{Twierdzenie}~[AN+NUM].

  \item \textbf{M5 --- most masa--sprz\k{e}\.zenie
    (prop.~\ref{prop:alphas-phiFP}):}
    $\alpha_s(M_Z) = N_c^3 g_0^{\,e}/(8\PhiZero) = 0.1190$
    ($1.2\sigma$ od PDG). Ta sama sta\l{}a $g_0^{\,e}$
    determinuje stosunek mas leptonowych i~si\l{}\k{e}
    oddzia\l{}ywa\'n silnych. Status: \statuslabel{Propozycja}~[AN+NUM].

  \item \textbf{Sektor kosmologiczny z~poprawk\k{a} FRW:}
    Wsp\'o\l{}czynnik $3\dot\psi^2/\psi$ (z~$c_0\psi$) nie zmienia
    predykcji inflacyjnych ($n_s$, $r$) --- atraktor
    Starobinsky'ego jest odporny na struktur\k{e} kinetyczn\k{a}.
    $\kappa = 3/(4\PhiZero) = 0.0304$: wci\k{a}\.z poni\.zej LLR.
    Weryfikacja: \texttt{cosmo\_frw\_verification\_v47.py} (18/18 PASS).

  \item \textbf{Master test v47:}
    Skrypt \texttt{tgp\_master\_consistency\_v47.py}:
    \textbf{50/51~PASS}, 1~oczekiwany FAIL
    (Koide bez przesuni\k{e}cia dla kwark\'ow --- potwierdza R12).

  \item \textbf{O12 --- $\sin^2\theta_W = 3/(N_c^2{+}N_c{+}1) = 3/13$
    (prop.~\ref{prop:sin2thetaW}):}
    Zero-parametrowa formu\l{}a daje $0{,}23077$ vs PDG $0{,}23122$
    (odch.\ $-0{,}19\%$). Mianownik $1 + N_c + N_c^2 = 13$
    jest sum\k{a} geometryczn\k{a} wymiar\'ow kolorowych.
    Status O12: awans \textbf{Otwarty}~$\to$~\textbf{[AN+HYP]}.

  \item \textbf{Czynnik $1/\varphi$ w~rurze kolorowej
    (prop.~\ref{prop:X-phi-color-tube}):}
    $\mathcal{A} = a_\Gamma/\varphi$ motywowany przez mod
    wt\'orny $\varphi$-FP: rura kolorowa ,,wyczerpuje'' jeden
    tryb solitonu na wi\k{a}zanie. Awans M4 do~\textbf{[AN+NUM]}.

  \item \textbf{O13 --- anomalia ABJ z~dynamicznym~$\Phi$
    (prop.~\ref{prop:anomaly-Phi-robust}):}
    Hiperladunki s\k{a} topologicznymi sta\l{}ymi
    (niezale\.znymi od $\Phi$); tw.~Adlera--Bardeena
    gwarantuje dok\l{}adno\'s\'c jednoobiegow\k{a}.
    Status O13: awans \textbf{Otwarty}~$\to$~\textbf{[AN]}.

  \item \textbf{Argument $k \le 2$ dla $\sin^2\theta_W$:}
    Saturacja tensorowa w~$d = 3$ wyklucza $k{=}3$
    (brak wolnych stopni swobody substratowych).
    Awans O12 do~\textbf{[AN]}.

  \item \textbf{Napi\k{e}cie z~EHT (O9) --- cz\k{e}\'sciowo zamkni\k{e}te (v47b):}
    Soliton TGP nie posiada sfery fotonowej.
    Analiza ISCO + profilu jasno\'sci
    (\texttt{bh\_isco\_raytrace\_tgp.py}) pokaza\l{}a:
    TGP ISCO przy $R_{\rm phys} \approx 0{,}55\,r_s$,
    pier\'scie\'n emisyjny przy $b \approx 0{,}6\,r_s$ (max.\ ugięcie $76°$, 1~przej\'scie dysku).
    TGP \textbf{produkuje} pier\'scie\'n z~ISCO, ale mniejszy i~s\l{}abszy ni\.z w~OTW.
    Napi\k{e}cie \textbf{zmniejszone} --- rozr\'o\.znialno\'s\'c wymaga ngEHT.
    Status O9: \textbf{CZ\k{E}\'SCIOWO ZAMKNI\k{E}TY}.

  \item \textbf{Sektor neutrinowy (v47b):}
    Predykcje: normal ordering, $\Sigma m_\nu \approx 62{,}9$~meV
    ($< 72$~meV Planck+DESI), $m_3/m_2 \approx 5{,}4$.
    Uwaga: ``$K(\nu) = 1/2$'' nie odpowiada standardowemu
    parametrowi Koide'a ($Q_K \in [1,3]$) --- wymaga sformalizowania
    w~kontek\'scie mechanizmu see-saw (Majorana).
    5~nowych test\'ow dodanych do master consistency.

  \item \textbf{Tabela zunifikowana (v47b):}
    \texttt{tgp\_unified\_predictions\_v47.py}:
    \textbf{24 predykcje z~2 parametr\'ow} (w tym 2 nowe z~rury kolorowej + QCD-corrected $\sin^2\theta_W$),
    22/24 w~$1\sigma$, 23/24 w~$2\sigma$.
    \texttt{tgp\_master\_consistency\_v47.py}:
    \textbf{55/56~PASS}, 1~oczekiwany FAIL
    (w tym 5 nowych test\'ow H1--H5: Koide-ODE).

  \item \textbf{Poprawka QCD do $\sin^2\theta_W$ (v47b):}
    Formu\l{}a~\eqref{eq:sin2-corrected}: $\sin^2\theta_W
    = N_c/(1+N_c+N_c^2 - b_0\alpha_s/N_c)$ z~wsp\'o\l{}czynnikiem
    beta QCD $b_0 = (11N_c{-}2N_f)/(12\pi)$.
    Redukcja odchylenia: \textbf{$11{,}3\sigma \to 0{,}6\sigma$}.
    Status O12: awans do~\textbf{[AN]}.

  \item \textbf{Rura kolorowa --- wariacyjne napi\k{e}cie strunowe (v47b):}
    Poprawiony potencja\l{} energetyczny:
    $V_E(\varphi) = \varphi^8/8 - \varphi^7/7 + 1/56$,
    z~$V_E(1) = 0$, $V_E'' > 0$ (pr\'o\.znia jest minimum).
    Gaussowski ansatz daje $\hat\sigma = \pi A^2$ (perturbacyjnie)
    z~$A_{\rm color} = C_F\alpha_s/(\pi\PhiZero) \approx 2 \times 10^{-3}$.

    \textbf{Kluczowy wynik:}
    $\mathcal{A} = a_\Gamma/\varphi = C_F^2\,\alpha_s^2(M_Z)$
    z~odchyleniem \textbf{0{,}04\%} przy $\alpha_s = 0{,}1179$ (PDG).
    Status M4/R12: awans do~\textbf{[AN+NUM]}.
    Skrypty: \texttt{color\_tube\_variational\_tgp.py},
    \texttt{color\_tube\_ode\_tgp.py}.
    Formalne propozycje: prop.~\ref{prop:X-tube-string-tension},
    prop.~\ref{prop:X-A-from-tube-tension}.
    7~nowych test\'ow (G1--G7) w~master consistency.

  \item \textbf{O-L5 --- Koide z~ODE (v47b, sek.~\ref{sec:T-koide-ode-v47b}):}
    Numeryczna weryfikacja \l{}a\'ncucha
    $d{=}3 \to k{=}4 \to N_{\rm gen}{=}3 \to \varphi\text{-FP}
    \to A_{\rm tail}(\text{ODE}) \to Q_K = 3/2$.
    Relacja Koide'go \textbf{nie jest} uniwersaln\k{a} w\l{}asno\'sci\k{a} ODE
    (losowe tr\'ojki: $Q_K \in [1{,}01;\,2{,}40]$).
    Warunek Koide'go w~$(p,q)$ daje r\'ownanie kwadratowe
    z~rozwi\k{a}zaniem $q_1 = 4{,}1010$, zgodnym z~ODE na $0{,}001\%$.
    Nowa relacja wirialna: $T_{\rm kin}/C_{\rm cross} \approx 1$
    z~dok\l{}adno\'sci\k{a} $0{,}2\%$ dla trzech generacji.
    \.Zadna z~8 testowanych form aproksymacyjnych nie odtwarza $Q_K = 3/2$
    --- wymaga pe\l{}nego rozwi\k{a}zania ODE.
    Status O-L5: \textbf{CZ\k{E}\'SCIOWO ZAMKNI\k{E}TY [NUM]}.
    Otwarty krok: dow\'od analityczny $A_{\rm tail} \to Q_K = 3/2$.
\end{enumerate}
\end{remark}

\begin{remark}[Odpowied\'z na audyt zewn\k{e}trzny (v47b)]%
\label{rem:audit-response}
Audyt z~2026-04-09 zidentyfikowa\l{} pi\k{e}\'c g\l{}\'ownych luk.
Poni\.zej podsumowanie dzia\l{}a\'n podj\k{e}tych w~sesji v47b.

\begin{enumerate}[nosep]
  \item \textbf{Brak mostu $\Gamma \to \Phi$ (CG-1, CG-3)}:
    Dodano dwie propozycje do dod.~Q
    (prop.~\ref{prop:Q-weak-continuum}:
    s\l{}abe twierdzenie continuum $H^1_{\rm loc}$;
    prop.~\ref{prop:Q-alpha-from-phi-squared}:
    $\alpha_{\rm eff} = 2$ z~$\Phi = \phi^2$).
    Status: \textbf{[SZKIC]}, nie pe\l{}ne domkni\k{e}cie.

  \item \textbf{Bilans predykcyjny --- mieszanie kalibracji z~predykcjami}:
    Utworzono uczciw\k{a} 4-klasow\k{a} taksonomie:
    \texttt{tgp\_prediction\_taxonomy\_v47.py}.
    Wynik: \textbf{2 wej\'scia [I]},
    \textbf{14 predykcji zerowych [D]} (ratio~7:1),
    5~odzysk\'ow [R], 4~prospektywne [P].
    17/19 testowalnych w~$1\sigma$.

  \item \textbf{Brak s\l{}ownika formulacji}:
    Dodano tabele przej\'s\'c Formulacja~A/B
    (uw.~\ref{rem:formulation-dictionary} w~sek.~\ref{sec:formalizm}).
    Zasada: Formulacja~A kanoniczna od sesji v33.

  \item \textbf{Brak mechanicznego uzasadnienia $m_0$}:
    Sformalizowano $\mathcal{A} = m_0\,m_1/m_3$ jako jedyny
    naturalny parametr rury kolorowej
    (prop.~\ref{prop:X-A-from-binding}),
    ze zgodno\'sci\k{a} $0{,}3\%$ mi\k{e}dzy sektorami.
    \textbf{Uzupe\l{}nienie v47b:} wariacyjne napi\k{e}cie strunowe
    z~ansatzem gaussowskim daje
    $\mathcal{A} = C_F^2\,\alpha_s^2$ z~odchyleniem $0{,}04\%$
    (prop.~\ref{prop:X-A-from-tube-tension}).
    Luka \textbf{istotnie domkni\k{e}ta}.

  \item \textbf{Sektor cechowania jako ,,plugin''}:
    Uwaga audytu s\l{}uszna, ale ju\.z oznaczona
    jako [Program] w~mapie statusu.
    Zasada defektowej hierarchii
    (amplituda $\to$ faza $\to$ multiplet)
    odnotowana jako kierunek przysz\l{}ego programu.
\end{enumerate}

\medskip\noindent
\textbf{Uwaga krytyczna wobec audytu:}
Audyt niedoszacowa\l{} si\l{}y predykcyjnej TGP.
Twierdzi\l{}, \.ze ,,wiele PASS to kalibracje lub to\.zsamo\'sci''.
4-klasowa taksonomia pokazuje, \.ze \textbf{14 niezale\.znych
predykcji z~2 parametr\'ow} to ratio 7:1 --- nietrywialne
nawet po ostrym filtrze.
Ponadto \l{}a\'ncuch
$d{=}3 \to k{=}4 \to N_{\rm gen}{=}3 \to Q_K{=}3/2 \to r^2{=}N{-}1$
jest zamkni\k{e}ty analitycznie (sesja v47) i~nie by\l{} dost\k{e}pny
audytorowi.
\end{remark}
